\subsection{Sweep-level comparison}

Across the 150 successful runs (75 hybrid, 75 SVM), the hybrid model outperformed the SVM baseline on patient-level discrimination metrics over most of the sweep landscape. The performance advantage was visible when averaging over window sizes (Figure~\ref{fig:sweep-window-roc}) and when averaging over \texttt{inner-k} settings (Figure~\ref{fig:sweep-innerk-roc}). The best individual hybrid run occurred at 6 s windows with \texttt{inner-k=1}, whereas the best SVM run occurred at 8 s windows with \texttt{inner-k=30}.

To assess statistical robustness across the structured sweep, we applied three correction tiers to the permutation p-values. Treating the sweep as two model-family comparisons (the primary independent factor) yields Bonferroni-corrected $p_\text{corr}=0.016$, confirming that the result survives at $\alpha=0.05$. A max-T permutation correction (10,000 iterations with shared label permutations across all 150 configurations, accounting for inter-configuration correlations) yields $p_\text{corr}=0.302$; the null maximum distribution has mean 0.768 and 95th percentile 0.908, placing the observed maximum of 0.816 within the range expected under random label permutations. The full-sweep Bonferroni ($p_\text{corr}=1.00$) is mechanically uninformative given the 150-test multiplier and is reported for completeness. We adopt the structured 2-family correction as the primary corrected p-value because the window-size and inner-k factors represent exploratory levels of a single hypothesis rather than independent comparisons. Across the sweep, 2 of 150 runs (1.3\%) yielded raw $p<0.05$, and 88\% of runs had $p\ge 0.50$, consistent with a null-dominated landscape from which the best hybrid configuration meaningfully separates.

\begin{paperinlinefigure}
  \centering
  \includegraphics[width=0.94\linewidth]{sweep_window_roc_auc.png}
  \caption{Patient-level ROC-AUC by window size across the sweep. Stars mark the best individual hybrid and SVM runs.}
  \label{fig:sweep-window-roc}
\end{paperinlinefigure}

Figure~\ref{fig:sweep-innerk-roc} shows that the relationship between ROC-AUC and \texttt{inner-k} differed across the two model families in the low-budget region. The hybrid model retained clearly stronger ROC-AUC across the sparsest budgets and achieved the highest individual ROC-AUC in the entire sweep at \texttt{inner-k=1}. The linear SVM did not show the same pattern and reached its best ROC-AUC only at the broader \texttt{inner-k=30} setting. Sparse per-fold selection therefore aligned with peak discrimination only for the more expressive model.

\begin{paperinlinefigure}
  \centering
  \includegraphics[width=0.94\linewidth]{sweep_inner_k_roc_auc.png}
  \caption{Patient-level ROC-AUC by \texttt{inner-k} across the sweep. The hybrid model retained the strongest discrimination in the sparsest budget region.}
  \label{fig:sweep-innerk-roc}
\end{paperinlinefigure}

Figures~\ref{fig:sweep-jaccard} and \ref{fig:sweep-unique-features} show a second related pattern. Increasing \texttt{inner-k} increased hard-set overlap and the total number of unique selected features. Larger feature budgets did not produce the strongest patient-level discrimination. The sweep concentrated its best hybrid performance at the sparsest end of the per-fold selection range. This pattern suggests that model capacity and per-fold sparsity mattered together rather than independently.

\begin{paperinlinefigure}
  \centering
  \includegraphics[width=0.94\linewidth]{sweep_jaccard_inner_k.png}
  \caption{Mean pairwise Jaccard overlap by \texttt{inner-k}. Hard-set overlap increased with larger per-fold feature budgets.}
  \label{fig:sweep-jaccard}
\end{paperinlinefigure}

\begin{paperinlinefigure}
  \centering
  \includegraphics[width=0.94\linewidth]{sweep_unique_features_inner_k.png}
  \caption{Unique selected features by \texttt{inner-k}. Larger per-fold budgets produced markedly larger discovered feature pools.}
  \label{fig:sweep-unique-features}
\end{paperinlinefigure}

\subsection{Best-run patient-level performance}

The best hybrid configuration achieved patient-level ROC-AUC 0.816 [95\% CI: 0.600--0.990], PR-AUC 0.579, balanced accuracy 0.786 [95\% CI: 0.571--0.967], accuracy 0.810 [95\% CI: 0.619--0.952], F1 0.714 [95\% CI: 0.364--0.941], and MCC 0.571 over 21 held-out participants (Figure~\ref{fig:main-roc}; Table~\ref{tab:best-run-comparison}). The permutation test yielded a nominal raw $p=0.008$ based on 500 iterations. Correcting for the structured sweep (2 model families) gives $p_\text{corr}=0.016$, which survives at $\alpha=0.05$; under the more conservative max-T correction across all 150 configurations, $p_\text{corr}=0.302$ (Table~\ref{tab:best-run-comparison} footnote). Bootstrap confidence intervals are wide, reflecting the small sample size of 21 participants. The PR-AUC of 0.579 indicates that the model retained useful precision-recall behavior in this cohort.

\begin{paperinlinefigure}
  \centering
  \includegraphics[width=0.94\linewidth]{main_patient_roc_curve.png}
  \caption{Patient-level ROC curve for the best hybrid run. The shaded region shows the 95\% bootstrap confidence band.}
  \label{fig:main-roc}
\end{paperinlinefigure}

The corresponding participant-level confusion counts were 12 true negatives, 5 true positives, 2 false positives, and 2 false negatives. This error profile is consistent with the balanced-accuracy summary and shows that the strongest run retained useful sensitivity without sacrificing specificity. Figure~\ref{fig:main-metrics} shows the same run in condensed metric form with 95\% bootstrap confidence intervals.

The best SVM baseline reached balanced accuracy 0.643 and ROC-AUC 0.510 [95\% CI: 0.119--0.881], with a non-significant permutation result (raw $p=0.469$, uncorrected). This weaker performance came with a much larger per-fold feature budget (\texttt{inner-k=30}) and a far larger overall candidate set. The contrast indicates that the sparse hybrid pipeline extracted stronger patient-level discrimination from a much smaller subset of features.

\begin{paperinlinefigure}
  \centering
  \includegraphics[width=0.94\linewidth]{main_patient_metric_summary.png}
  \caption{Patient-level metric summary for the best hybrid run. Error bars show 95\% bootstrap confidence intervals.}
  \label{fig:main-metrics}
\end{paperinlinefigure}

\begin{paperwidetable}
  \centering
  \caption{Best-performing configuration for each model family. Bootstrap confidence intervals are from 1,000 resamples at the 95\% level.}
  \label{tab:best-run-comparison}
  \resizebox{\linewidth}{!}{\begin{tabular}{lcccccccccc}
\toprule
Model & Window & Inner-k & ROC-AUC & PR-AUC & Bal.\ Acc. & F1 & MCC & Feat./fold & Unique feat. & Mean Jaccard \\
\midrule
Hybrid CNN-LSTM & 6 & 1 & 0.816 & 0.579 & 0.786 & 0.714 & 0.571 & 1.0 & 7 & 0.190 \\
Linear SVM & 8 & 30 & 0.510 & 0.411 & 0.643 & 0.533 & 0.277 & 30.0 & 103 & 0.360 \\
\bottomrule
\end{tabular}
}

  \smallskip
  \footnotesize Bootstrap 95\% CIs from 1,000 resamples. Hybrid permutation test: raw $p=0.008$; primary correction (2 model families) $p_\text{corr}=0.016$; robustness checks: max-T across 150 configurations $p_\text{corr}=0.302$, full-sweep Bonferroni $p_\text{corr}=1.00$.
\end{paperwidetable}

\subsection{Intermediate baselines and ablation at the sparsest budget}

To determine whether the hybrid architecture's advantage at \texttt{inner-k=1} reflected generic benefits of non-linearity or specific architectural properties, we evaluated eight classical classifiers and two neural baselines under the identical configuration (6 s windows, \texttt{inner-k=1}, SMOTE, LOPO equalization, same dataset and folds). All ten baselines performed at or below chance (Table~\ref{tab:intermediate-baselines}). The best classical model (Random Forest) reached ROC-AUC 0.296, compared with 0.816 for the hybrid. A 2-layer MLP (ROC-AUC 0.204) and a GAP variant of the hybrid architecture that removes self-attention and feature-pyramid fusion (ROC-AUC 0.235) both failed to exceed chance. No baseline achieved a significant permutation p-value (all $p \ge 0.9$).

An additional ablation removed SMOTE from the best hybrid configuration. Performance fell from ROC-AUC 0.816 to 0.582 (balanced accuracy 0.500, F1 0.308), indicating that SMOTE contributes substantially to the hybrid's performance in the 1-dimensional feature regime — consistent with the observation that synthetic oversampling provides the only mechanism for exposing class-boundary information when a single scalar feature is available.

These findings indicate that the hybrid CNN-LSTM's performance at extreme sparsity is not attributable to non-linearity alone, nor to SMOTE alone. The convolutional, recurrent, attention, and feature-pyramid stages, together with SMOTE-based training, appear jointly necessary to extract discriminatory signal from a single feature per fold.

\begin{paperwidetable}
  \centering
  \caption{Intermediate baseline classifiers and ablations evaluated at the best hybrid configuration (6~s, \texttt{inner-k=1}). All models used identical data, feature selection, and validation splits. Neural baselines and ablation trained on NVIDIA RTX 3090.}
  \label{tab:intermediate-baselines}
  \footnotesize
  \begin{tabular}{lcccccc}
\toprule
Model & ROC-AUC & Bal.\ Acc. & MCC & F1 & $p$ (perm.) \\
\midrule
Hybrid CNN-LSTM (best) & \textbf{0.816} & \textbf{0.786} & \textbf{0.571} & \textbf{0.714} & 0.008 \\
\midrule
Random Forest & 0.296 & 0.321 & $-$0.337 & 0.222 & 0.930 \\
Extra Trees & 0.286 & 0.500 & 0.000 & 0.375 & 0.934 \\
AdaBoost & 0.286 & 0.393 & $-$0.204 & 0.250 & 0.942 \\
Gradient Boosting & 0.276 & 0.393 & $-$0.204 & 0.250 & 0.936 \\
Decision Tree & 0.276 & 0.500 & 0.000 & 0.308 & 0.946 \\
KNN & 0.255 & 0.464 & $-$0.067 & 0.353 & 0.974 \\
SVM-RBF & 0.204 & 0.429 & $-$0.135 & 0.333 & 0.978 \\
SVM-Linear & 0.204 & 0.357 & $-$0.277 & 0.300 & 0.988 \\
\bottomrule
\end{tabular}

\smallskip
\footnotesize All models evaluated with the same configuration: window size 6~s, \texttt{inner-k}=1, SelectKBest ($F$-classif), SMOTE enabled, LOPO group equalization enabled, patient-level mean aggregation. Permutation p-values from 500 iterations. The hybrid result is shown for reference (not part of this sweep).

\end{paperwidetable}

The winning hybrid run selected an average of 1 feature per fold and yielded only 7 unique features that appeared in at least one fold across the entire participant-level validation process (Table~\ref{tab:hybrid-biomarkers}). That degree of sparsity is one of the most distinctive findings in the sweep. Hard-set overlap was modest, with a mean pairwise Jaccard value of 0.190 and a mean Kuncheva index of 0.187. Figure~\ref{fig:biomarker-frequency} shows the small candidate pool, while the overlap statistics indicate that recurrence was concentrated in a few repeatedly reappearing features rather than in a broadly shared foldwise set.

A small set of features emerged repeatedly, though not in every fold, and the selected features were dominated by temporal and frontal-temporal descriptors. The most frequently selected feature was \texttt{tp10\_spectral\_entropy}, selected in 8 of 21 folds. The next most frequent were \texttt{left\_frontal\_temporal\_diff\_gamma}, \texttt{left\_frontal\_temporal\_diff\_beta}, and \texttt{af7\_zero\_crossings}. We note that three of the seven features listed in Table~\ref{tab:hybrid-biomarkers} appear in only 1 of 21 folds; the term ``recurrent'' should be interpreted with this sparsity in mind, and the sign-consistency values for these single-fold features carry limited interpretive weight.

\begin{paperinlinefigure}
  \centering
  \includegraphics[width=0.94\linewidth]{biomarker_selection_frequency.png}
  \caption{Selection-frequency ranking for the features identified from the best hybrid run.}
  \label{fig:biomarker-frequency}
\end{paperinlinefigure}

\begin{paperwidetable}
  \centering
  \caption{Features from the best hybrid run that appeared in at least one fold. $\Delta$ selection is remission minus non-remission selection share. Sign consistency is 1.000 for features appearing in multiple folds; single-fold entries should be interpreted with caution.}
  \label{tab:hybrid-biomarkers}
  \resizebox{\linewidth}{!}{\begin{tabular}{lcccccc}
\toprule
Feature & Folds & Share & $\Delta$ sel. & Sign cons. & Mean $d$ & Effect direction \\
\midrule
tp10\_spectral\_entropy & 8 & 0.381 & -0.571 & 1.000 & -0.935 & Non-remission > remission \\
left\_frontal\_temporal\_diff\_gamma & 4 & 0.190 & -0.071 & 1.000 & -1.059 & Non-remission > remission \\
left\_frontal\_temporal\_diff\_beta & 3 & 0.143 & 0.000 & 1.000 & -0.915 & Non-remission > remission \\
af7\_zero\_crossings & 3 & 0.143 & 0.429 & 1.000 & 0.814 & Remission > non-remission \\
tp10\_hurst & 1 & 0.048 & 0.143 & 1.000 & -1.010 & Non-remission > remission \\
af7\_tp10\_diff\_gamma & 1 & 0.048 & 0.143 & 1.000 & -1.203 & Non-remission > remission \\
frontal\_asymmetry\_ratio\_delta & 1 & 0.048 & -0.071 & 1.000 & 0.794 & Remission > non-remission \\
\bottomrule
\end{tabular}
}
\end{paperwidetable}

\subsection{Effect-direction consistency and class-conditional patterns}

The recurrent hybrid biomarkers showed perfect mean sign consistency, with all recurrent features retaining the same effect direction across the folds in which they were selected. \texttt{tp10\_spectral\_entropy} and the frontal-temporal gamma and beta difference terms were repeatedly associated with the non-remission group, whereas \texttt{af7\_zero\_crossings} was more frequently selected in remission folds. The class-conditional selection frequencies make that separation visible at the level of foldwise recurrence (Figure~\ref{fig:biomarker-class-conditional}).

\begin{paperinlinefigure}
  \centering
  \includegraphics[width=0.94\linewidth]{figure6_class_conditional_selection.png}
  \caption{Class-conditional selection frequencies for the recurrent biomarkers from the best hybrid run. Remission is shown in red and non-remission in blue.}
  \label{fig:biomarker-class-conditional}
\end{paperinlinefigure}

The effect-size summaries followed the same broad pattern, with \texttt{af7\_zero\_crossings} retaining a positive mean effect size and the main TP10 and frontal-temporal features retaining negative mean effects (Figure~\ref{fig:biomarker-effect-size}). This distinction matters clinically. The selection-frequency deltas do not justify definitive mechanistic claims. They provide a concrete interpretation target for follow-up validation. Taken together, these results indicate that useful patient-level discrimination can be achieved with an extremely sparse feature budget, while the recurrent biomarkers retain interpretable directionality even when fold-to-fold set overlap remains modest.

\subsection{Aggregation robustness}

Patient-level metrics in the main analysis were derived by averaging window probabilities within each participant. To assess sensitivity to outlier windows, we recomputed the best hybrid run's patient-level metrics using median aggregation and 25\% trimmed-mean aggregation. Median aggregation yielded ROC-AUC 0.857, PR-AUC 0.620, and MCC 0.571; trimmed-mean aggregation yielded ROC-AUC 0.837 and MCC 0.571. Both alternative aggregation methods matched or exceeded the original mean-aggregated performance (ROC-AUC 0.816), confirming that the main results are not driven by outlier windows and that the reported mean-based metrics are conservative rather than inflated.

\begin{paperinlinefigure}
  \centering
  \includegraphics[width=0.94\linewidth]{figure7_effect_size.png}
  \caption{Effect-size stability for the recurrent biomarkers from the best hybrid run. Positive values indicate remission-skewed effects and negative values indicate non-remission-skewed effects.}
  \label{fig:biomarker-effect-size}
\end{paperinlinefigure}
